% ===================================================================
% LH-05c: Lehrerhinweise - Das Verb gustar
% Abgeleitet aus LE-05c (Verlauf, Differenzierung, Fehler)
% ===================================================================

\documentclass[11pt, a4paper]{article}

\usepackage[utf8]{inputenc}
\usepackage[T1]{fontenc}
\usepackage[ngerman]{babel}
\usepackage{helvet}
\renewcommand{\familydefault}{\sfdefault}

\usepackage[margin=2cm]{geometry}
\usepackage{enumitem}
\usepackage{tabularx}
\usepackage{booktabs}
\usepackage{longtable}

\usepackage{xcolor}
\usepackage{tcolorbox}

\usepackage{fancyhdr}
\usepackage{lastpage}
\usepackage{amssymb}

% Farben
\definecolor{spanisch}{HTML}{c41e3a}
\definecolor{grau}{HTML}{888888}
\definecolor{hellgrau}{HTML}{f5f5f5}
\definecolor{basis}{HTML}{2a7a4b}
\definecolor{standard}{HTML}{2c5aa0}
\definecolor{erweiterung}{HTML}{7b1fa2}
\definecolor{loesung}{HTML}{c62828}

\newcommand{\fachfarbe}{spanisch}

% Variablen
\newcommand{\thema}{Das Verb gustar (gefallen/mögen)}
\newcommand{\sequenz}{05}
\newcommand{\block}{c}
\newcommand{\fach}{Spanisch}
\newcommand{\klasse}{7}
\newcommand{\dauer}{45 Min. (Einzelstunde)}

% Kopf-/Fußzeile
\pagestyle{fancy}
\fancyhf{}
\renewcommand{\headrulewidth}{0.4pt}
\renewcommand{\footrulewidth}{0.4pt}
\lhead{\fach{} -- Klasse \klasse}
\chead{\textbf{LH-\sequenz\block{} (Lehrerhinweise)}}
\rhead{}
\lfoot{}
\rfoot{Seite \thepage\ von \pageref{LastPage}}

% Differenzierungsmarker
\newcommand{\B}{\textcolor{basis}{\textbf{[B]}}}
\newcommand{\Std}{\textcolor{standard}{\textbf{[S]}}}
\newcommand{\E}{\textcolor{erweiterung}{\textbf{[E]}}}
\newcommand{\loesung}[1]{\textcolor{loesung}{\textbf{#1}}}
\newcommand{\es}[1]{\textcolor{\fachfarbe}{\textbf{#1}}}

% Boxen
\newtcolorbox{kurzplan}{
    colback=hellgrau,
    colframe=\fachfarbe,
    boxrule=1pt,
    arc=0pt,
    fonttitle=\bfseries,
    title=Kurzplan
}

\newtcolorbox{differenzierung}{
    colback=white,
    colframe=grau,
    boxrule=0.5pt,
    arc=0pt,
    fonttitle=\bfseries,
    title=Differenzierung
}

\newtcolorbox{fehlerbox}{
    colback=white,
    colframe=loesung,
    boxrule=0.5pt,
    arc=0pt,
    fonttitle=\bfseries,
    title=Häufige Fehler
}

\newtcolorbox{materialbox}{
    colback=white,
    colframe=\fachfarbe,
    boxrule=0.5pt,
    arc=0pt,
    fonttitle=\bfseries,
    title=Material-Checkliste
}

% ===================================================================
\begin{document}

\begin{center}
    {\Large\bfseries\textcolor{\fachfarbe}{Lehrerhinweise: \thema}}\\[0.3em]
    {\normalsize LH-\sequenz\block{} | \dauer{} | Sequenz 5: Mi tiempo libre}
\end{center}

\vspace{10pt}

% ===================================================================
% 1. KURZPLAN
% ===================================================================

\section*{1. Stundenverlauf (Kurzplan)}

\begin{tabularx}{\textwidth}{|c|l|X|l|}
\hline
\textbf{Zeit} & \textbf{Phase} & \textbf{Inhalt} & \textbf{Material} \\
\hline
0--5 & Einstieg & Bildimpuls: ``Me gusta el fútbol. ¿Y a ti?'' & Bilder \\
\hline
5--10 & Aktivierung & Beispiele sammeln: gusta vs. gustan -- Was fällt auf? & Tafel \\
\hline
10--20 & Input & gustar-Konstruktion: Pronomen + gusta/gustan & Tafel, AB \\
\hline
20--28 & Übung 1 & EA: Zuordnung gusta/gustan (AB Aufg. 1-2) & AB-05c \\
\hline
28--38 & Übung 2 & PA: Dialoge ``¿Qué te gusta?'' & Bildkarten \\
\hline
38--43 & Sicherung & 2-3 Paare präsentieren & -- \\
\hline
43--45 & Abschluss & Merkkasten sichern, HA & -- \\
\hline
\end{tabularx}

% ===================================================================
% 2. DIFFERENZIERUNG
% ===================================================================

\section*{2. Differenzierung}

\begin{differenzierung}
\textbf{Legende:} \B{} Basis (BR) | \Std{} Standard (MR) | \E{} Erweiterung (GY)

\vspace{0.5em}
\textbf{WICHTIG:} Diese Marker sind NUR für die Lehrkraft!
\end{differenzierung}

\vspace{1em}

\begin{tabularx}{\textwidth}{|l|X|X|X|}
\hline
& \B{} \textbf{Basis} & \Std{} \textbf{Standard} & \E{} \textbf{Erweiterung} \\
\hline
\textbf{Aufg. 1} & Mit Wortbank (gusta/gustan) & Ohne Hilfe & + Begründung \\
\hline
\textbf{Aufg. 2} & Pronomen vorgegeben & Pronomen selbst wählen & + weitere Sätze \\
\hline
\textbf{Aufg. 3} & Dialog mit Vorlage & Dialog mit Stichworten & Freier Dialog \\
\hline
\textbf{Aufg. 4} & 3 Sätze (me gusta) & 5 Sätze (verschiedene) & + Negation, alle Pronomen \\
\hline
\end{tabularx}

\vspace{1em}

\textbf{Konkrete Hinweise:}
\begin{itemize}
    \item \B{} LRS-SuS: Zeitzugabe (+10 Min.), Pronomen-Tabelle als Hilfe
    \item \B{} DaZ-SuS: Kontrastierung Deutsch-Spanisch betonen
    \item \Std{} Standardgruppe: Partnerarbeit mit wechselnden Partnern
    \item \E{} Leistungsstarke: Alle Pronomen verwenden, Negation einbauen
\end{itemize}

% ===================================================================
% 3. HÄUFIGE FEHLER
% ===================================================================

\section*{3. Häufige Fehler}

\begin{fehlerbox}
\begin{tabularx}{\textwidth}{|l|X|X|}
\hline
\textbf{Fehler} & \textbf{Beispiel} & \textbf{Reaktion} \\
\hline
Interferenz ``Ich mag'' & *``Yo gusto el fútbol'' & Kontrastierung: ``Mir gefällt'' -- Visualisierung \\
\hline
gusta/gustan verwechselt & *``Me gusta los deportes'' & Faustregel: Was gefällt? Singular/Plural? \\
\hline
Pronomen vergessen & *``Gusta la música'' & Betonen: IMMER Pronomen (me, te, le...) \\
\hline
Artikel fehlt & *``Me gusta fútbol'' & Hinweis: ``Me gusta \textbf{el} fútbol'' \\
\hline
Verwechslung me/mi & *``Mi gusta...'' & me = mir (vor Verb), mi = mein (vor Nomen) \\
\hline
\end{tabularx}
\end{fehlerbox}

% ===================================================================
% 4. MATERIAL-CHECKLISTE
% ===================================================================

\section*{4. Material-Checkliste}

\begin{materialbox}
\begin{itemize}[label=$\square$]
    \item AB-05c.pdf (24 Kopien)
    \item ML-05c.pdf (1x für Lehrkraft)
    \item Lehrwerk: Qué pasa 1, S. 73-74
    \item Bildkarten Freizeitaktivitäten (12 Sets)
    \item Tafel/Whiteboard für Visualisierung
    \item Optional: LP-05c.html (Tablets/PCs)
\end{itemize}
\end{materialbox}

% ===================================================================
% 5. LÖSUNGEN
% ===================================================================

\section*{5. Lösungen AB-05c}

\textbf{Merkkasten:}
\begin{itemize}
    \item Pronomen: me, te, le, nos, os, les
    \item Beispiele: Te gusta..., Le gusta..., Nos gusta..., Os gusta..., Les gusta...
\end{itemize}

\vspace{0.5em}

\textbf{Aufgabe 1:} (je 1P = 5P)
\begin{enumerate}[label=\alph*)]
    \item \loesung{gusta} (Singular: el fútbol)
    \item \loesung{gustan} (Plural: los deportes)
    \item \loesung{gusta} (Infinitiv: leer)
    \item \loesung{gustan} (Plural: las películas)
    \item \loesung{gusta} (Singular: la música)
\end{enumerate}

\textbf{Aufgabe 2:} (je 1P = 4P)
\begin{enumerate}[label=\alph*)]
    \item \loesung{Te gusta}
    \item \loesung{le gusta}
    \item \loesung{Nos gustan}
    \item \loesung{Os gusta}
\end{enumerate}

\textbf{Aufgabe 3:} (max. 5P)
\begin{itemize}
    \item A: ¿Qué te gusta hacer?
    \item B: Me gusta jugar al fútbol. / Me gusta [Aktivität].
    \item A: ¿Te gustan los videojuegos? / ¿Te gusta [...]?
    \item B: Sí, me gustan. / No, no me gustan.
\end{itemize}

\textbf{Aufgabe 4:} (max. 6P)
\begin{itemize}
    \item Korrekte gustar-Konstruktion: 3P
    \item Richtige Wahl gusta/gustan: 2P
    \item Verschiedene Pronomen/Negation: 1P
\end{itemize}

% ===================================================================
% 6. HAUSAUFGABE
% ===================================================================

\section*{6. Hausaufgabe}

\begin{itemize}
    \item AB-05c: Aufgabe 4 fertigstellen
    \item Vokabeln lernen: gustar-Konstruktion, Freizeitaktivitäten
    \item Optional: Lernpfad LP-05c.html bearbeiten
\end{itemize}

% ===================================================================
% 7. REFLEXIONSNOTIZEN
% ===================================================================

\section*{7. Reflexionsnotizen (nach Durchführung)}

\begin{tabularx}{\textwidth}{|l|X|}
\hline
\textbf{Was lief gut?} & \\[2em]
\hline
\textbf{Was lief nicht?} & \\[2em]
\hline
\textbf{Zeitmanagement} & \\[2em]
\hline
\textbf{Für nächstes Mal} & \\[2em]
\hline
\end{tabularx}

\end{document}
